% bundle_geometry.tex
%
% How to recompute an ar_track_alvar bundle file (bundles/room.xml) from
% observations, derived from scratch.  Written to be re-read, not skimmed:
% every step is spelled out, including the ones that look obvious and are
% exactly where sign errors live.
%
% Build with:  pdflatex bundle_geometry.tex

\documentclass[11pt,a4paper]{article}

\usepackage[margin=2.4cm]{geometry}
\usepackage{amsmath,amssymb}
\usepackage{booktabs}
\usepackage{xcolor}
\usepackage[colorlinks=true,linkcolor=blue!50!black,urlcolor=blue!50!black]{hyperref}
\usepackage{enumitem}

\newcommand{\T}[2]{T_{#1 \leftarrow #2}}   % transform: takes #2-coords, gives #1-coords
\newcommand{\R}[2]{R_{#1 \leftarrow #2}}
\newcommand{\vt}[1]{\mathbf{#1}}
\newcommand{\xh}{\hat{\vt{x}}}
\newcommand{\yh}{\hat{\vt{y}}}
\newcommand{\zh}{\hat{\vt{z}}}

\title{Recomputing an AR bundle file from scratch\\
       \large the geometry, with nothing skipped}
\date{}

\begin{document}
\maketitle

\section*{What this document is for}

A bundle file (\texttt{bundles/room.xml}) tells \texttt{ar\_track\_alvar} where several
printed markers sit \emph{relative to each other}. If the markers get moved, the file is
wrong, and every pose that depends on it is wrong in a way that still looks plausible.

The file cannot be measured conveniently with a tape measure (the numbers are corner
positions in a frame that may not physically exist), but it \emph{can} be computed from
two observations. This document derives that computation from first principles, so that
it can be redone after any rearrangement without rediscovering the algebra.

Nothing here is harder than vectors and matrix multiplication. The difficulty is
entirely bookkeeping: which frame is a number expressed in, and which direction does a
transform go. Sections~\ref{sec:frames}--\ref{sec:inverse} are about exactly that, and
they are the part worth reading slowly.

\tableofcontents

% ======================================================================
\section{Points, frames, and what a ``pose'' is}
\label{sec:frames}

A \emph{point} in the room is a physical thing: a corner of a printed square. It has no
numbers of its own. Numbers appear only once a \emph{frame} is chosen.

A frame is an origin plus three perpendicular unit axes $\xh,\yh,\zh$, forming a
right-handed set ($\xh \times \yh = \zh$). The coordinates of a point in that frame are
the three distances you travel along those axes to reach it.

The same physical point therefore has different coordinates in different frames. Writing
$\vt{p}$ with no indication of frame is the single most common source of error in this
kind of work, so in this document every quantity carries its frame.

Four frames appear:

\begin{center}
\begin{tabular}{ll}
\toprule
$C$ & the camera (\texttt{radius\_infra1\_optical\_frame}) \\
$M$ & the \emph{master} frame of the bundle (marker index 0; in our room, fictional) \\
$11$ & the frame of the printed marker with id 11 \\
$12$ & the frame of the printed marker with id 12 \\
\bottomrule
\end{tabular}
\end{center}

A \emph{pose} of frame $B$ relative to frame $A$ is two things together: where $B$'s
origin sits, expressed in $A$ coordinates, and how $B$'s axes are turned relative to
$A$'s. That is exactly the information needed to convert coordinates from $B$ into $A$,
which is why ``pose of $B$ in $A$'' and ``transform from $B$ to $A$'' are the same object.

% ======================================================================
\section{Rotation matrices: what the columns mean}

A rotation is stored as a $3\times3$ matrix $\R{A}{B}$. The definition worth memorising
is this one:

\begin{quote}
\textbf{The columns of $\R{A}{B}$ are the axes of frame $B$, written in $A$ coordinates.}
\end{quote}

So if $B$'s $\xh$ axis points along $A$'s $\zh$ axis, the first column of $\R{A}{B}$ is
$(0,0,1)^{\!\top}$. That is the whole content of the matrix; everything else follows.

Why that definition makes the matrix do the right thing: a point with $B$-coordinates
$\vt{p}_B = (p_1,p_2,p_3)^{\!\top}$ is, by the meaning of coordinates, the combination
$p_1 \xh_B + p_2 \yh_B + p_3 \zh_B$. Writing each of those axis vectors in $A$
coordinates and adding them up is precisely the matrix--vector product:
\begin{equation}
  \vt{p}_A \;=\; p_1\,(\text{col}_1) + p_2\,(\text{col}_2) + p_3\,(\text{col}_3)
           \;=\; \R{A}{B}\, \vt{p}_B .
\end{equation}

\paragraph{Properties, and how to check a matrix is really a rotation.}
Because the columns are perpendicular unit vectors:
\begin{align}
  \R{A}{B}^{\!\top} \R{A}{B} &= I  &&\text{(columns orthonormal)}\\
  \det \R{A}{B} &= +1            &&\text{(right-handed, not mirrored)}\\
  \R{A}{B}^{-1} &= \R{A}{B}^{\!\top} &&\text{(the inverse is free)}
\end{align}
The determinant check matters. A matrix with $\det = -1$ is a reflection: it will look
almost right, it will pass the orthonormality test, and it will produce a world that is
mirror-imaged. If a fit ever returns $\det = -1$, something upstream is wrong --- do not
``fix'' it by flipping a sign at random.

\paragraph{How big is a rotation?} The angle $\theta$ of a rotation matrix $R$ is
\begin{equation}
  \theta = \arccos\!\left( \frac{\operatorname{trace}(R) - 1}{2} \right).
  \label{eq:angle}
\end{equation}
This is useful for sanity checks: ``how far did my correction turn the marker?'' is
answered by applying~\eqref{eq:angle} to the correction matrix. Clamp the argument into
$[-1,1]$ before calling \texttt{arccos}, because rounding will occasionally hand you
$1.0000000002$.

% ======================================================================
\section{Homogeneous transforms: rotation and translation in one object}

Rotation alone cannot express ``and also, the origin is over there''. The full transform
is
\begin{equation}
  \vt{p}_A = \R{A}{B}\,\vt{p}_B + \vt{t}_{A\leftarrow B},
\end{equation}
where $\vt{t}_{A\leftarrow B}$ is the position of $B$'s origin in $A$ coordinates.

Carrying a matrix and a vector around separately is error-prone, so both are packed into
one $4\times4$ matrix:
\begin{equation}
  \T{A}{B} =
  \begin{pmatrix}
    \R{A}{B} & \vt{t}_{A\leftarrow B} \\[2pt]
    \vt{0}^{\!\top} & 1
  \end{pmatrix},
  \qquad
  \begin{pmatrix} \vt{p}_A \\ 1 \end{pmatrix}
  = \T{A}{B}
  \begin{pmatrix} \vt{p}_B \\ 1 \end{pmatrix}.
\end{equation}
The trailing $1$ on the point is what lets the translation column act; the bottom row
$(0,0,0,1)$ keeps that $1$ in place after multiplication. There is no deeper meaning ---
it is a bookkeeping trick that makes composition work.

\subsection{Composition: the cancellation rule}
\label{sec:compose}

The reason the arrow notation is worth the extra typing:
\begin{equation}
  \boxed{\;\T{A}{B}\,\T{B}{C} = \T{A}{C}\;}
  \label{eq:chain}
\end{equation}
The inner labels must match and they cancel, exactly like units. Read it right to left:
start with $C$-coordinates, the right-hand factor turns them into $B$-coordinates, the
left-hand factor turns those into $A$-coordinates.

If a product you have written does not cancel --- say $\T{A}{B}\,\T{C}{B}$ --- it is
meaningless, and no amount of transposing will rescue it. One of the two factors needs
inverting first. Checking cancellation before computing anything catches most mistakes
in this subject.

\subsection{The inverse}
\label{sec:inverse}

Reversing a transform is not ``put a minus sign on everything''. Starting from
$\vt{p}_A = R\,\vt{p}_B + \vt{t}$ and solving for $\vt{p}_B$:
\begin{align}
  \vt{p}_A - \vt{t} &= R\,\vt{p}_B \\
  R^{\!\top}(\vt{p}_A - \vt{t}) &= \vt{p}_B
    \qquad\text{(using } R^{-1} = R^{\!\top}\text{)}\\
  \vt{p}_B &= R^{\!\top}\vt{p}_A \;+\; \bigl(-R^{\!\top}\vt{t}\bigr).
\end{align}
So the inverse has rotation $R^{\!\top}$ and translation $-R^{\!\top}\vt{t}$:
\begin{equation}
  \T{A}{B}^{-1} \;=\; \T{B}{A} \;=\;
  \begin{pmatrix}
    R^{\!\top} & -R^{\!\top}\vt{t} \\[2pt]
    \vt{0}^{\!\top} & 1
  \end{pmatrix}.
  \label{eq:inverse}
\end{equation}
The $-R^{\!\top}\vt{t}$ term is the one people drop. Geometrically: the new translation
has to be the old origin's position \emph{seen from the other frame}, which involves
rotating it, not just negating it.

% ======================================================================
\section{What is actually in the bundle file}

Here is the relevant part of the file, before the correction:

\begin{quote}\ttfamily\small
\begin{tabbing}
\texttt{<marker index="0" status="1">}\\
\quad\=\texttt{<corner x="6.5"~~y="-6.5"~~~~z="0" />}\\
\>\texttt{<corner x="6.5"~~y="6.5"~~~~~z="0" />}\\
\>\texttt{<corner x="-6.5" y="6.5"~~~~~z="0" />}\\
\>\texttt{<corner x="-6.5" y="-6.5"~~~~z="0" />}\\
\texttt{</marker>}\\[4pt]
\texttt{<marker index="12" status="1">}\\
\>\texttt{<corner x="-59"~~y="-4"~~z="51" />}\\
\>\texttt{<corner x="-59"~~y="9"~~~z="51" />}\\
\>\texttt{<corner x="-59"~~y="9"~~~z="64" />}\\
\>\texttt{<corner x="-59"~~y="-4"~~z="64" />}\\
\texttt{</marker>}
\end{tabbing}
\end{quote}

Three facts, none of which the file states:

\begin{enumerate}[leftmargin=*]
\item \textbf{The units are centimetres.} The markers are $13\,$cm squares and the
      half-width $6.5$ appears everywhere. Everything published on ROS topics is in
      metres. Convert once, early, and write down which unit you are in.

\item \textbf{All corner coordinates are expressed in the master frame $M$.} The master
      is the marker with \texttt{index="0"}. This is why its own corners are the clean
      symmetric numbers $(\pm 6.5, \pm 6.5, 0)$: a marker expressed in its own frame
      sits at the origin, flat in the $z=0$ plane.

\item \textbf{The master need not exist physically.} In our room, marker 0 was never
      printed or placed. That is allowed --- $M$ is just the coordinate system the other
      markers are described in. It has a consequence we return to in
      Section~\ref{sec:gauge}.
\end{enumerate}

\paragraph{The local corner convention.} Reading marker 0's entry tells us the order and
orientation alvar expects for a marker's own corners. With half-width $s = 6.5$\,cm:
\begin{equation}
  \vt{l}_0 = \begin{pmatrix}s\\-s\\0\end{pmatrix}\!,\quad
  \vt{l}_1 = \begin{pmatrix}s\\s\\0\end{pmatrix}\!,\quad
  \vt{l}_2 = \begin{pmatrix}-s\\s\\0\end{pmatrix}\!,\quad
  \vt{l}_3 = \begin{pmatrix}-s\\-s\\0\end{pmatrix}.
  \label{eq:local}
\end{equation}
So in a marker's own frame: $\xh$ runs from corner 3 to corner 0, $\yh$ runs from corner
0 to corner 1, and $\zh$ points out of the printed face. \emph{The order of the four
corners in the XML is part of the data.} Permuting them silently rotates or flips the
marker.

% ======================================================================
\section{From four corners to a transform}
\label{sec:corners2T}

Given a marker's four corners $\vt{w}_0 \ldots \vt{w}_3$ in master coordinates, we want
$\T{M}{k}$, the pose of that marker in the master frame.

Use the definition from Section 2 --- the columns of a rotation matrix are the child
frame's axes expressed in the parent. We know from~\eqref{eq:local} which pairs of
corners define those axes, so we just read them off:

\begin{align}
  \xh &= \frac{\vt{w}_0 - \vt{w}_3}{\lVert \vt{w}_0 - \vt{w}_3 \rVert}
        &&\text{(the local $+x$ direction)} \label{eq:xax}\\[2pt]
  \yh &= \frac{\vt{w}_1 - \vt{w}_0}{\lVert \vt{w}_1 - \vt{w}_0 \rVert}
        &&\text{(the local $+y$ direction)} \label{eq:yax}\\[2pt]
  \zh &= \xh \times \yh
        &&\text{(out of the face, right-handed)} \label{eq:zax}
\end{align}
and
\begin{equation}
  \R{M}{k} = \bigl(\, \xh \;\big|\; \yh \;\big|\; \zh \,\bigr),
  \qquad
  \vt{t}_{M\leftarrow k} = \tfrac{1}{4}\textstyle\sum_{i=0}^{3} \vt{w}_i .
  \label{eq:fit}
\end{equation}

The translation is just the centroid of the four corners, because the local corners
in~\eqref{eq:local} are symmetric about the origin and therefore average to zero.

\paragraph{Why $\zh$ comes from a cross product rather than the corners.} The four
corners are coplanar, so they contain no information about which side of the plane is
``front'' beyond the handedness convention. Taking $\xh \times \yh$ enforces
$\det = +1$ by construction.

\paragraph{Checks worth doing every time.}
\begin{itemize}[leftmargin=*,itemsep=1pt]
\item $\lVert \vt{w}_0 - \vt{w}_3 \rVert$ and $\lVert \vt{w}_1 - \vt{w}_0 \rVert$ should
      both equal the marker width ($13$\,cm). If they differ, the XML entry is not a
      square and was hand-fitted.
\item $\xh \cdot \yh$ should be $0$. If it is not, the entry is not even a rectangle.
\item Applying~\eqref{eq:fit} to marker 0's own corners must give the identity. It is a
      free unit test of your implementation, and it confirms you have the corner order
      right.
\end{itemize}

\paragraph{When the corners are not perfectly consistent.} Marker 11's entry in our file
has side lengths $13.0$ and $12.9$\,cm, so no rigid transform reproduces it exactly.
Equations~\eqref{eq:xax}--\eqref{eq:zax} will still return something, but it quietly
favours whichever two corners you happened to use. The honest alternative is the
\emph{Kabsch} fit, which uses all four corners and minimises the total squared error:

\begin{enumerate}[leftmargin=*,itemsep=1pt]
\item Centre both sets: $\tilde{\vt{l}}_i = \vt{l}_i - \bar{\vt{l}}$ and
      $\tilde{\vt{w}}_i = \vt{w}_i - \bar{\vt{w}}$.
\item Build $H = \sum_i \tilde{\vt{l}}_i \tilde{\vt{w}}_i^{\!\top}$ (a $3\times3$ matrix).
\item Take the singular value decomposition $H = U S V^{\!\top}$.
\item Set $R = V \operatorname{diag}(1,1,d) U^{\!\top}$ with
      $d = \operatorname{sign}\bigl(\det(V U^{\!\top})\bigr)$.
\item Then $\vt{t} = \bar{\vt{w}} - R\,\bar{\vt{l}}$.
\end{enumerate}
The $d$ term is what prevents the answer from coming out as a reflection. The leftover
per-corner distances after applying $R,\vt{t}$ are a genuine residual: for a clean entry
they are $0$, and for marker 11 they are $0.5$\,mm at every corner.

% ======================================================================
\section{From a transform back to four corners}

To write a marker back into the XML, run Section~\ref{sec:corners2T} backwards. For each
of the four local corners in~\eqref{eq:local}:
\begin{equation}
  \vt{w}_i \;=\; \R{M}{k}\,\vt{l}_i \;+\; \vt{t}_{M \leftarrow k},
  \qquad i = 0,1,2,3,
  \label{eq:back}
\end{equation}
keeping the same order, and converting metres back to centimetres.

Do not expect the new numbers to be as tidy as the old ones. The original entry had
every $x$ equal to $-59$ because the marker was assumed perfectly parallel to a wall.
A measured marker is not, so all four $x$ values will differ slightly. That is the
measurement, not a mistake.

% ======================================================================
\section{What the camera gives us}

\texttt{ar\_track\_alvar} publishes visualisation markers on
\texttt{visualization\_marker}. Two namespaces matter:

\begin{center}
\begin{tabular}{ll}
\toprule
\texttt{basic\_shapes}, id $=k$ & where alvar thinks \emph{the individual marker $k$} is \\
\texttt{main\_shapes}, id $=0$  & where alvar thinks \emph{the bundle's master frame} is \\
\bottomrule
\end{tabular}
\end{center}

Both are poses in the camera frame $C$. The second one is the interesting one, and the
reason is worth stating carefully:

\begin{quote}
To report the master pose, alvar must combine what it \emph{sees} (a marker) with what it
has been \emph{told} (that marker's position in the bundle file). So
\texttt{main\_shapes} is part measurement, part assumption. \textbf{If the bundle file is
wrong, \texttt{main\_shapes} is wrong --- and it is wrong by a different amount depending
on which marker is currently visible.} That difference is the signal we exploit.
\end{quote}

Hence the experiment:

\begin{enumerate}[leftmargin=*,itemsep=1pt]
\item Cover marker 12. Only 11 is visible. Record \texttt{main\_shapes}. Call it $A$.
\item Cover marker 11. Only 12 is visible. Record \texttt{main\_shapes}. Call it $B$.
\item \textbf{Do not move the camera between the two recordings.}
\end{enumerate}

If the bundle file were correct, $A$ and $B$ would be identical. They are not.

% ======================================================================
\section{The derivation}
\label{sec:derivation}

Name the quantities, each with its frames:

\begin{center}
\begin{tabular}{lll}
\toprule
symbol & meaning & known? \\
\midrule
$A = \T{C}{M}^{(11)}$ & master pose, computed by alvar from marker 11 & measured \\
$B = \T{C}{M}^{(12)}$ & master pose, computed by alvar from marker 12 & measured \\
$X = \T{M}{12}$       & marker 12's pose in the master frame, \emph{per the old XML} & known, wrong \\
$C_{12} = \T{C}{12}$  & where marker 12 truly is, seen by the camera & unknown \\
$X_{\text{true}}$     & what the XML \emph{should} say & what we want \\
\bottomrule
\end{tabular}
\end{center}

\paragraph{Step 1: write down what alvar did to produce $B$.}
It measured marker 12 directly, giving $\T{C}{12}$, then chained through the bundle file
to reach the master:
\begin{equation}
  B = \T{C}{M} = \underbrace{\T{C}{12}}_{\text{measured}} \; \underbrace{\T{12}{M}}_{\text{from XML}}
    = C_{12} \, X^{-1}.
  \label{eq:Bdef}
\end{equation}
Check the cancellation: $C \leftarrow 12$ then $12 \leftarrow M$ gives $C \leftarrow M$. Good.
Note $\T{12}{M} = X^{-1}$, because the XML stores $X = \T{M}{12}$, the other direction.

\paragraph{Step 2: eliminate the unknown.}
Rearranging~\eqref{eq:Bdef} gives the camera's actual view of marker 12:
\begin{equation}
  C_{12} = B\,X.
  \label{eq:C12}
\end{equation}
This is the key move. $C_{12}$ is a \emph{measurement} that was never published directly,
but it can be reconstructed, because alvar's error came entirely from the XML and we
know what the XML said.

\paragraph{Step 3: decide what the master frame means.}
The correct XML entry is marker 12's pose in the master frame:
\begin{equation}
  X_{\text{true}} = \T{M}{12} = \T{M}{C}\;\T{C}{12} = A^{-1}\,C_{12},
\end{equation}
where we took $A$ --- the estimate from marker 11 --- as the definition of $\T{C}{M}$.
Section~\ref{sec:gauge} explains why that is a legitimate choice and not an assumption
being smuggled in.

\paragraph{Step 4: substitute~\eqref{eq:C12} and finish.}
\begin{equation}
  \boxed{\;X_{\text{true}} \;=\; A^{-1}\,B\,X\;}
  \label{eq:result}
\end{equation}

Read the result right to left: take the old XML entry, apply the discrepancy $A^{-1}B$
between the two master estimates, and you have the corrected entry. The factor $A^{-1}B$
is \emph{exactly} the error that was in the file --- it is the transform that moves
``where 12 claimed the master was'' onto ``where 11 says the master is''.

\paragraph{Sanity check on the formula.} If the XML had been right, $A$ and $B$ would be
equal, so $A^{-1}B = I$ and $X_{\text{true}} = X$: the formula leaves a correct file
alone. That check costs nothing and catches an inverted factor immediately.

% ======================================================================
\section{Turning the published numbers into matrices}

ROS publishes orientations as quaternions $(q_x, q_y, q_z, q_w)$ --- note that ROS puts
$w$ \emph{last}, while most textbooks put it first. Getting this wrong produces a
rotation that is wrong but not obviously wrong.

First normalise, $\lVert q \rVert = 1$, then
\begin{equation}
  R =
  \begin{pmatrix}
    1 - 2(q_y^2 + q_z^2) & 2(q_x q_y - q_z q_w) & 2(q_x q_z + q_y q_w) \\
    2(q_x q_y + q_z q_w) & 1 - 2(q_x^2 + q_z^2) & 2(q_y q_z - q_x q_w) \\
    2(q_x q_z - q_y q_w) & 2(q_y q_z + q_x q_w) & 1 - 2(q_x^2 + q_y^2)
  \end{pmatrix}.
  \label{eq:q2R}
\end{equation}
Always verify $R^{\!\top}R = I$ and $\det R = +1$ afterwards. If the quaternion was not
normalised, the first check fails and tells you so.

\subsection{Averaging many samples}

A single frame is noisy. Record a few seconds and average.

\paragraph{Translations} average normally: add them up, divide by the count.

\paragraph{Rotations do not.} Averaging rotation matrices entry by entry gives a matrix
that is not a rotation. There is also a sign trap: $q$ and $-q$ describe the
\emph{same} rotation, so a naive mean of samples that happen to differ in sign tends
toward zero.

The standard fix:
\begin{enumerate}[leftmargin=*,itemsep=1pt]
\item Normalise every quaternion.
\item Pick one as a reference. For every other sample, if $q \cdot q_{\text{ref}} < 0$,
      replace $q$ by $-q$. Now they all live in the same hemisphere.
\item Stack them as the rows of an $N \times 4$ matrix $Q$.
\item Form $Q^{\!\top}Q$ (a $4\times4$ symmetric matrix) and take its eigenvector for the
      largest eigenvalue. That eigenvector is the average quaternion.
\end{enumerate}

Step 4 looks like magic but is just least squares: it finds the unit quaternion most
aligned with all the samples at once. For tightly clustered samples it agrees with the
naive mean to several decimals; it is the sign handling in step 2 that does most of the
work in practice.

\paragraph{Report the spread, not just the mean.} The standard deviation of the
translations and the maximum angle between each sample and the mean tell you whether the
recording was any good. For the data this document is based on, the 58 samples of $B$
had a translation spread of $0.3$\,mm and a rotation spread of $0.03^\circ$ --- i.e.\
the camera and markers really were static.

Note what that does and does not prove: it is \emph{precision}, not accuracy. A
systematically wrong pose can be extremely repeatable.

% ======================================================================
\section{Worked example with the real numbers}

Units are metres throughout this section; conversion back to centimetres happens at the
very end.

\subsection{Input $A$ --- master seen through marker 11}

One recorded sample:
\begin{equation*}
  \vt{t}_A = \begin{pmatrix} 0.181993 \\ 0.633671 \\ -0.018187 \end{pmatrix},
  \qquad
  q_A = (0.236931,\; 0.694806,\; -0.650365,\; 0.195279).
\end{equation*}
Through~\eqref{eq:q2R}:
\begin{equation*}
  R_A =
  \begin{pmatrix}
    -0.811460 &  0.583247 & -0.036821 \\
     0.075237 &  0.041779 & -0.996290 \\
    -0.579545 & -0.811220 & -0.077783
  \end{pmatrix}.
\end{equation*}

\subsection{Input $B$ --- master seen through marker 12}

Average of 58 samples:
\begin{equation*}
  \vt{t}_B = \begin{pmatrix} -0.213286 \\ 0.584412 \\ 0.262373 \end{pmatrix},
  \qquad
  q_B = (0.223091,\; 0.678418,\; -0.676997,\; 0.177917),
\end{equation*}
\begin{equation*}
  R_B =
  \begin{pmatrix}
    -0.837152 &  0.543597 & -0.060659 \\
     0.061799 & -0.016188 & -0.997957 \\
    -0.543468 & -0.839190 & -0.020042
  \end{pmatrix}.
\end{equation*}

\subsection{Input $X$ --- the old XML entry for marker 12}

Corners $(-59,-4,51),\,(-59,9,51),\,(-59,9,64),\,(-59,-4,64)$ in cm, i.e.
$(-0.59,-0.04,0.51)$ etc.\ in metres. Applying~\eqref{eq:xax}--\eqref{eq:fit}:

\begin{align*}
  \xh &= \frac{(-0.59,-0.04,0.51) - (-0.59,-0.04,0.64)}{0.13}
       = \frac{(0,0,-0.13)}{0.13} = (0,0,-1), \\
  \yh &= \frac{(-0.59,0.09,0.51) - (-0.59,-0.04,0.51)}{0.13}
       = \frac{(0,0.13,0)}{0.13} = (0,1,0), \\
  \zh &= \xh \times \yh = (0,0,-1) \times (0,1,0) = (1,0,0),
\end{align*}
so, stacking those as \emph{columns},
\begin{equation*}
  X =
  \begin{pmatrix}
     0 & 0 & 1 & -0.590 \\
     0 & 1 & 0 &  0.025 \\
    -1 & 0 & 0 &  0.575 \\
     0 & 0 & 0 &  1
  \end{pmatrix}.
\end{equation*}
The centroid $(-0.590,\,0.025,\,0.575)$ is the average of the four corners, as promised.

It is worth pausing on what this matrix says: marker 12's own $x$ axis points along the
master's $-z$, its $y$ along the master's $+y$, and its face normal along the master's
$+x$. In words, it is on a wall to one side, facing inward --- which matches the physical
room, and is a good sign that the corner order was read correctly.

\subsection{Step 1: invert $A$}

Using~\eqref{eq:inverse}, $R_A^{\!\top}$ and $-R_A^{\!\top}\vt{t}_A$:
\begin{equation*}
  A^{-1} =
  \begin{pmatrix}
    -0.811460 &  0.075237 & -0.579545 &  0.089465 \\
     0.583247 &  0.041779 & -0.811220 & -0.147374 \\
    -0.036821 & -0.996290 & -0.077783 &  0.636607 \\
     0        &  0        &  0        &  1
  \end{pmatrix}.
\end{equation*}
Observe that the rotation block is $R_A$ transposed --- rows and columns swapped --- while
the translation is \emph{not} $-\vt{t}_A$. That is~\eqref{eq:inverse} in action.

\subsection{Step 2: the discrepancy $A^{-1}B$}

\begin{equation*}
  A^{-1}B =
  \begin{pmatrix}
     0.998929 &  0.044023 & -0.014245 &  0.154450 \\
    -0.044812 &  0.997143 & -0.060814 & -0.460199 \\
     0.011527 &  0.061388 &  0.998047 &  0.041808 \\
     0        &  0        &  0        &  1
  \end{pmatrix}.
\end{equation*}

This single matrix is the entire error that was sitting in the bundle file. Decoded:
\begin{itemize}[leftmargin=*,itemsep=1pt]
\item translation $(15.4,\,-46.0,\,4.2)$\,cm --- the two estimates of the master frame
      disagreed by about half a metre;
\item rotation angle, from~\eqref{eq:angle} with
      $\operatorname{trace} = 2.994119$: $\theta = 4.39^\circ$.
\end{itemize}
A near-identity rotation block with a large translation is the signature of ``a marker
got moved'', as opposed to ``a marker got re-oriented''. Worth noting for diagnosis next
time.

\subsection{Step 3: apply it}

\begin{equation*}
  X_{\text{true}} = (A^{-1}B)\,X =
  \begin{pmatrix}
     0.014245 & 0.044023 &  0.998929 & -0.442008 \\
     0.060814 & 0.997143 & -0.044812 & -0.443800 \\
    -0.998047 & 0.061388 &  0.011527 &  0.610419 \\
     0        & 0        &  0        &  1
  \end{pmatrix}.
\end{equation*}
Centre moved from $(-59.0,\,2.5,\,57.5)$ to $(-44.2,\,-44.4,\,61.0)$\,cm, a shift of
$49.3$\,cm, with the marker turned by $4.39^\circ$.

\subsection{Step 4: back to corners}

Via~\eqref{eq:back} with $s = 0.065$\,m, then $\times 100$:

\begin{quote}\ttfamily\small
\begin{tabbing}
\texttt{<marker index="12" status="1">}\\
\quad\=\texttt{<corner x="-44.39" y="-50.47" z="54.16" />}\\
\>\texttt{<corner x="-43.82" y="-37.50" z="54.95" />}\\
\>\texttt{<corner x="-44.01" y="-38.29" z="67.93" />}\\
\>\texttt{<corner x="-44.58" y="-51.26" z="67.13" />}\\
\texttt{</marker>}
\end{tabbing}
\end{quote}

The four $x$ values are no longer identical, and the $z$ values no longer come in two
matching pairs. That is the $4.39^\circ$ showing up in the numbers.

% ======================================================================
\section{Why the master frame is a free choice}
\label{sec:gauge}

In Step 3 of Section~\ref{sec:derivation} we declared that $A$ --- the estimate through
marker 11 --- \emph{defines} the master frame. Why is that allowed, when $A$ is just as
much a measurement as $B$?

Because the master frame is not observable. Marker 0 was never placed. Nothing in the
room, in the camera, or in the data knows where $M$ ``really'' is. What the bundle file
genuinely encodes is the \emph{relative} geometry of markers 11 and 12; the master is
only the coordinate system that relationship is written in.

Shifting $M$ anywhere changes every entry in the file, but changes no prediction the
system makes --- the camera poses that come out the other end are identical. A degree of
freedom with no observable consequence is called a \emph{gauge} freedom, and the usual
practice is to pin it to something convenient. Leaving marker 11's entry untouched is
convenient: it means a re-measurement only ever edits the markers that moved.

Two consequences:
\begin{itemize}[leftmargin=*,itemsep=1pt]
\item Marker 11's entry has become \emph{definitional}. It cannot be ``wrong'' any more
      in the sense that matters, because it is what $M$ means. (It can still be wrong in
      the sense of not being a $13\,$cm square --- see Section~\ref{sec:corners2T}.)
\item If you ever want $M$ to coincide with something physical --- the floor, say, as was
      originally intended --- that is a separate decision requiring a separate
      measurement, and it changes every entry in the file at once.
\end{itemize}

% ======================================================================
\section{What this method cannot tell you}

\paragraph{There is no residual.} Two pose observations, one rigid transform to solve
for, six unknowns and six constraints: the problem is \emph{exactly determined}. It fits
perfectly, always, by construction. Applying the correction makes $A$ and $B$ agree
because that is literally the equation we solved --- not because the answer is right.

So this procedure cannot report its own accuracy. To get a residual you need redundancy:
more viewpoints than unknowns. Record the pair from three or four camera positions and
solve for the one $X$ that best fits all of them; the leftover disagreement is then a
real error estimate.

\paragraph{The output inherits the weaker input.} The result is only as good as the worse
of $A$ and $B$. In the example, $B$ averaged 58 samples while $A$ was a single frame ---
so the final numbers carry single-frame noise regardless of how carefully $B$ was
measured. Record both with the same care.

\paragraph{Single-marker pose has an ambiguity.} Estimating a pose from one planar square
is ill-conditioned out of plane: there are two poses that project to nearly the same
image, differing by a flip about an axis in the marker plane. Solvers sometimes pick the
wrong one, especially for markers that are small in the image, viewed close to
head-on, or far away. A sudden tens-of-degrees change between consecutive samples is the
signature. Both of $A$ and $B$ are single-marker estimates, so this is the failure mode
to watch for.

\paragraph{A camera that moved invalidates everything.} The derivation assumes $A$ and
$B$ share one camera frame $C$. Nothing in the data reveals a bump between recordings.
Take both recordings back to back.

% ======================================================================
\section{The recipe, next time the room changes}

\begin{enumerate}[leftmargin=*]
\item Point the camera so that \emph{both} markers are comfortably in view, reasonably
      large in the image, and not edge-on. Then do not touch it.

\item Cover marker 12. Record \texttt{visualization\_marker} for a few seconds. Keep only
      the messages with \texttt{ns: "main\_shapes"}. This is $A$.

\item Uncover 12, cover 11. Record again, same way. This is $B$.

\item Average each set: mean translation, eigenvector-method quaternion. Print the
      spreads and look at them before continuing.

\item Build $A$ and $B$ as $4\times4$ matrices using~\eqref{eq:q2R}.

\item Read the \emph{old} XML entry for marker 12 and build $X$
      with~\eqref{eq:xax}--\eqref{eq:fit}. Convert cm to m. Check the side lengths are
      $13$\,cm and that marker 0's entry gives the identity.

\item Compute $X_{\text{true}} = A^{-1} B X$ using~\eqref{eq:inverse} for the inverse.

\item Report the shift and the rotation angle. A shift of the size you remember moving
      the marker is the single best confirmation available.

\item Generate the new corners with~\eqref{eq:back}, convert back to cm, write them into
      the file keeping the corner order.

\item \textbf{Verify independently}: run with both markers visible and watch
      \texttt{main\_shapes} while covering one, then the other. If the master pose does
      not jump, the relative geometry is right. This is the only check that is not
      circular.
\end{enumerate}

\paragraph{If there are more than two markers.} Repeat steps 2--9 once per extra marker,
always against the same reference $A$, and always leaving the reference marker's entry
alone. With $n$ markers you need $n$ recordings, one per marker, all from the same
camera pose.

\end{document}
